\subsection{The gap across models}\label{sec:models-res}
The picture does not depend on the detector (Table~\ref{tab:headline}; full matrices for every model are in the Supplementary Material). The four content models reach in-domain AUC between \InAUCmin{} and \InAUCmax{} and off-diagonal AUC between \OffAUCmin{} and \OffAUCmax{}, and all of them show the same split between ranking and operating-point loss (Table~\ref{tab:decomp}). Averaged word embeddings (\textsc{glove-lr}) are weaker in-domain (\GloveInAUC) and transfer no better (\GloveOffAUC).

The style model is the most revealing. With only 22 surface statistics it reaches an average in-domain AUC of \StyleInAUC{} (\StyleISOT{} on ISOT, \StyleWEL{} on WELFake), and its transfer AUC (\StyleOffAUC) is only a few points below that of the lexical models. A large part of what these benchmarks reward is form: length, casing, punctuation and attribution habits (Section~\ref{sec:shortcuts}).

Pretrained encoders do not change the conclusion. On identical 1{,}500-item subsets, frozen RoBERTa features give in-domain AUC \RobInAUC{} and off-diagonal AUC \RobOffAUC{}, against \SubLRInAUC{} and \SubLROffAUC{} for \textsc{tfidf-lr} refitted on the same splits. The paired off-diagonal difference is \RobMinusLRoff{} AUC (Wilcoxon $p=\RobMinusLRp$), and RoBERTa is better in \RobBetterCells{} of the 42 cells. Fine-tuning the upper four layers does not change this either. On the same 800 training and 300 test items per corpus (\FTseeds{} seeds), \textsc{roberta-ft} reaches in-domain AUC \FTInAUC{} and off-diagonal AUC \FTOffAUC{}, against \FTsubInAUC{} and \FTsubOffAUC{} for the paired \textsc{tfidf-lr}; the off-diagonal difference is \FTminusLRoff{} (Wilcoxon $p=\FTminusLRp$, better in \FTbetterCells{} of 42 cells). Its gap is again mostly ranking (\FTRankShare), EM again makes the operating point worse (\FTRecEM{} recovered), and rate transfer again recovers most of it (\FTRecRank). With only four layers trained on 800 items, the fine-tuned model is slightly weaker than \textsc{tfidf-lr} in-domain (\FTminusLRin{} AUC), so this is not the strongest transformer one could build; but its transfer matches that of the lexical model, and the shape of its gap is the same. This agrees with \citet{verma2026audit} and \citet{delloglio2026comparison}, who found that fully fine-tuned transformers gain in-domain but not, or not reliably, across corpora. Pairwise McNemar tests on the same target items (Holm-corrected over all comparisons) find significant differences from \textsc{tfidf-lr} in the following share of (cell, seed) comparisons: \McFracSig.
